\documentclass[12pt]{article}

\usepackage[margin=1in]{geometry}
\usepackage{booktabs}
\usepackage{array}
\usepackage{tabularx}
\usepackage{threeparttable}
\usepackage{hyperref}
\usepackage{setspace}
\usepackage{parskip}
\usepackage{amsmath}
\usepackage{amssymb}
\usepackage[T1]{fontenc}
\usepackage{lmodern}
\usepackage{microtype}

\hypersetup{
    colorlinks=true,
    linkcolor=blue,
    urlcolor=blue,
    citecolor=blue
}

\title{\textbf{Job Preferences and Wage Expectations\\Among Ghanaian University Graduates}}
\author{Remy Beauregard}
\date{May 2026}

\begin{document}

\maketitle

\begin{abstract}
\noindent
This report summarises preliminary findings from a baseline survey of 48 Ghanaian
university graduates on their wage expectations and the factors that influence their
job-choice decisions. Regression analysis reveals that professional development is positively associated
with both the minimum and maximum acceptable wage, while two stated social influences,
family and peer opinion, are associated only with the maximum: family opinion is
associated with a lower ceiling wage and peer opinion with a higher one. These are
exploratory patterns intended to inform the design of a pilot RCT.
\end{abstract}

\vspace{1em}
\hrule
\vspace{1.5em}

\setstretch{1.2}

%─────────────────────────────────────────────────────────────────────────────
\section{Background and Data}
%─────────────────────────────────────────────────────────────────────────────

The survey was administered to Ghanaian university graduates through KoboCollect.
Respondents provided their minimum and maximum willingness-to-accept (WTA) wages
(the floor and ceiling of pay they would consider for a job offer), along with
demographic information and a set of indicators for which factors (e.g.\ location,
hours flexibility, professional development) influence their employment decisions.

Because respondents entered wages in free text, responses arrived in diverse
formats: monthly figures (``GHC 2,000''), hourly rates (``15 Ghana Cedis per hour''),
daily rates, ranges (``1,500--2,000''), and even written-out words (``two thousand'').
All entries were algorithmically standardised to a \textbf{monthly GHC equivalent}
using a 40-hour, five-day working week (approximately 173 hours per month and
21.7 working days per month). Range entries were represented by their midpoint.

The analytical sample was constructed from the raw survey responses in two steps.
First, respondents who were not college graduates were dropped (grad\_code $= 3$;
$n = 1$ dropped). Second, respondents missing wage data for both the
minimum and maximum WTA, or missing any key demographic or job-factor variable, were
excluded ($n = 4$ dropped). The \textbf{baseline analytical sample
contains 48 respondents}. Summary statistics for key variables are shown in
Table~\ref{tab:sumstats}.

\begin{table}[h!]
\centering
\begin{threeparttable}
\caption{Summary Statistics ($n = 48$)}
\label{tab:sumstats}
\begin{tabular}{lcccccc}
\toprule
\textbf{Variable} & \textbf{Mean} & \textbf{SD} & $p_{5}$ & $p_{95}$ & \textbf{Min} & \textbf{Max} \\
\midrule
\multicolumn{7}{l}{\textit{Demographics}} \\[2pt]
Age &  27.5 &   3.6 &    24 &    34 &    24 &    43 \\
Female &  0.38 &  0.49 &  0.00 &  1.00 &     0 &     1 \\
Married &  0.10 &  0.31 &  0.00 &  1.00 &     0 &     1 \\
Has children &  0.08 &  0.28 &  0.00 &  1.00 &     0 &     1 \\[4pt]
\multicolumn{7}{l}{\textit{Wages (monthly GHC)}} \\[2pt]
Minimum WTA &   2248 &   1624 &    100 &   4000 &     18 &  10000 \\
Maximum WTA &   4260 &   2647 &    250 &  10000 &     30 &  10000 \\
\bottomrule
\end{tabular}
\begin{tablenotes}[flushleft]
\scriptsize
\item \textit{Notes:} Female, Married, and Has children are binary (0/1) indicators; their means equal the share of the sample with that characteristic. $p_5$ and $p_{95}$ for wage variables are the winsorization thresholds used in Table~\ref{tab:regsw}. Wages standardized to monthly GHC equivalents using a 40-hour, 5-day working week.
\end{tablenotes}
\end{threeparttable}
\end{table}


%─────────────────────────────────────────────────────────────────────────────
\section{Prevalence of Job-Choice Factors}
%─────────────────────────────────────────────────────────────────────────────

Respondents were asked which factors influence their job-choice decisions.
Table~\ref{tab:social} shows the reported prevalence of each factor.

\begin{table}[h!]
\centering
\caption{Respondents Citing Each Factor as Influencing Job Choice ($n = 48$)}
\label{tab:social}
\begin{tabular}{lcc}
\toprule
\textbf{Factor} & \textbf{Count} & \textbf{Share} \\
\midrule
Professional development & 37 & 77.1\% \\
Location                 & 36 & 75.0\% \\
Hours flexibility        & 36 & 75.0\% \\
Training                 & 30 & 62.5\% \\
Qualifications match     & 19 & 39.6\% \\
Promotion prospects      & 14 & 29.2\% \\
Family opinion           &  6 & 12.5\% \\
Peer opinion             &  1 &  2.1\% \\
\bottomrule
\multicolumn{3}{l}{\footnotesize \textit{Note:} The one respondent citing peer opinion is included among the six citing family opinion.}\\
\end{tabular}
\end{table}

The striking rarity of stated social influence, particularly peer opinion (cited by
only one respondent), sets up an important contrast with the regression results
discussed in Section~\ref{sec:reg}, where both social factors emerge as statistically
significant predictors of the \emph{maximum} acceptable wage.

%─────────────────────────────────────────────────────────────────────────────
\section{Regression Results}
\label{sec:reg}
%─────────────────────────────────────────────────────────────────────────────

We estimate four OLS regressions, each using heteroskedasticity-robust standard
errors. The outcome is either the minimum or the maximum monthly wage WTA
(in GHC). Predictors are eight binary job-factor indicators
(location, family opinion, peer opinion, hours flexibility, qualifications match,
promotion prospects, professional development, and training). Two specifications
add age and gender (female) as demographic controls.

Table~\ref{tab:regs} summarises the direction and significance of notable
coefficients. Full regression output is available in the Stata log. Four
respondents provided a minimum but not a maximum wage, so models (3) and (4)
use 44 observations.

\input{regtable_table.tex}

\subsection{Minimum Wage WTA}

Of all job-factor indicators, \textbf{professional development} is the only
significant predictor of a respondent's minimum acceptable wage (significant at
the 5\% level, $p < 0.05$, in both specifications). The positive coefficient
indicates that respondents who cite professional development as important in their
job choice name a \emph{higher} wage floor. This is consistent with a
self-selection story: individuals oriented toward career growth may have higher
reservation wages, either because they feel their human capital warrants more or
because they are willing to wait for a better match. The result is robust to
the inclusion of age and gender as controls.

No other job-factor indicator (location, hours flexibility, qualifications match,
promotion, or training) is significantly associated with the minimum WTA.

\subsection{Maximum Wage WTA}

Three job-factor indicators are significantly associated with the maximum
acceptable wage ($n = 44$):

\begin{itemize}
    \item \textbf{Professional development} is associated with a significantly
    \emph{higher} maximum WTA in both specifications ($p = 0.009$ without
    controls; $p = 0.006$ with controls). The coefficient is substantially
    larger for the ceiling than for the floor, suggesting career-oriented
    graduates set higher aspirations across the board.

    \item \textbf{Family opinion} is associated with a significantly \emph{lower}
    maximum WTA ($p = 0.001$ in both specifications). Respondents who report
    that family opinion matters in their job choice name a lower ceiling wage.
    One interpretation is that family networks may channel graduates toward
    sectors where salaries are typically modest, or may reflect conservative
    reference points for expected pay in the local graduate labour market.

    \item \textbf{Peer opinion} is associated with a significantly \emph{higher}
    maximum WTA ($p = 0.005$ without controls; $p = 0.026$ with controls).
    This could reflect peer networks that set high wage norms, or social
    comparison dynamics that raise wage aspirations.
\end{itemize}

The controlled specification (model 4) also reveals a significant negative
association between age and maximum WTA ($p = 0.044$), suggesting older
respondents in the sample set lower wage ceilings.

%─────────────────────────────────────────────────────────────────────────────
\section{Noteworthy Observations}
%─────────────────────────────────────────────────────────────────────────────

\paragraph{The stated--revealed disconnect.}
The most striking pattern in the data is the mismatch between how \emph{few}
respondents explicitly report social factors as important in their job decisions
(1 of 48 for peer opinion; 6 of 48 for family opinion) and the statistical
prominence of those same factors in the maximum WTA regressions, where both are
significant at the 1\% level. This disconnect may indicate that social influences
operate at a level respondents do not consciously attribute to their
decision-making, an interesting hypothesis to explore in subsequent work.

\paragraph{Asymmetry between social factors and professional development.}
Social factors (family and peer opinion) are significant only for the maximum
(ceiling) WTA, not the minimum (floor). Professional development, by contrast,
is significant for both: it raises the floor (by roughly GHC 860--870/month)
and raises the ceiling by a larger margin (roughly GHC 2,200--2,450/month).
This pattern suggests that career-oriented graduates set higher wages across the
board, while social reference points shape aspiration levels but not reservation
prices.

\paragraph{Salary excluded from regressions.}\label{sec:notes}
The salary job-factor indicator was not included as a predictor in any regression. This is appropriate: using ``does salary matter
to you?'' as a predictor of the actual salary you demand would introduce a form of
circularity (it measures the same underlying construct as the outcome), potentially
inflating coefficients and distorting inference on other factors.

\paragraph{Wage data heterogeneity.}
The free-text wage format produced responses in at least four different time units
(hourly, daily, monthly, ranges) and in multiple currencies and spellings. The
algorithmic standardisation procedure handles these cases, but some uncertainty
remains; for example, a respondent who wrote ``GHC 3,000 and above'' is treated
as having stated exactly GHC 3,000. Measurement error in the outcome variable
reduces statistical power and may affect regression estimates.

%─────────────────────────────────────────────────────────────────────────────
\section{Robustness: Winsorised Wages}
%─────────────────────────────────────────────────────────────────────────────

The raw wage data include extreme values at both ends (e.g.\ a minimum WTA of GHC~18
and a maximum of GHC~10,000 per month) that may reflect measurement or reporting
noise rather than genuine preferences. To assess whether these
extreme values drive the main results, we re-estimate all four regressions
after winsorising both wage outcomes at the 5th and 95th percentiles: values
below the 5th percentile are set to the 5th percentile value, and values above
the 95th are set to the 95th. This limits the influence of outliers without
discarding any observations.

Table~\ref{tab:regsw} reports the results. The winsorisation affected 4 minimum-wage
observations (2 below GHC~100 and 2 above GHC~4,000) and 2 maximum-wage observations
(below GHC~250; none exceeded the 95th percentile of GHC~10,000). All significant
findings are robust: professional development retains its sign, direction, and
significance for both the minimum ($p = 0.026$; $p = 0.035$ with controls) and maximum
WTA ($p = 0.009$; $p = 0.006$), family opinion remains negatively associated with the
maximum WTA at the 1\% level, and peer opinion remains positive and significant
($p = 0.005$ without controls; $p = 0.027$ with controls). Point estimates for
non-significant minimum-WTA predictors shift more (notably training, which roughly
halves, and hours flexibility, which moves near zero), though none cross a significance
threshold, suggesting those estimates are sensitive to a small number of extreme values.

\begin{table}[h!]
\centering
\begin{threeparttable}
\caption{OLS Regression Results: Winsorised WTA (5th/95th Percentile)}
\label{tab:regsw}
\begin{tabular}{l*{4}{>{\centering\arraybackslash}p{1.6cm}}}
\toprule
 & \multicolumn{2}{c}{\textbf{Min WTA}} & \multicolumn{2}{c}{\textbf{Max WTA}} \\
\cmidrule(lr){2-3}\cmidrule(lr){4-5}
\textbf{Job factor} & (1) & (2) & (3) & (4) \\
\midrule
\multicolumn{5}{l}{\textit{Significant predictors}} \\[2pt]
Professional development & $+$** & $+$** & $+$*** & $+$*** \\
Family opinion & ns & ns & $-$*** & $-$*** \\
Peer opinion & ns & ns & $+$*** & $+$** \\[4pt]
\multicolumn{5}{l}{\textit{Non-significant predictors}} \\[2pt]
Location & ns & ns & ns & ns \\
Hours flexibility & ns & ns & ns & ns \\
Qualifications match & ns & ns & ns & ns \\
Promotion prospects & ns & ns & ns & ns \\
Training & ns & ns & ns & ns \\
\midrule
Controls (age, female) & & \checkmark & & \checkmark \\
$N$ & 48 & 48 & 44 & 44 \\
\bottomrule
\end{tabular}
\begin{tablenotes}[flushleft]
\footnotesize
\item \textit{Note:} $+$ / $-$ indicates coefficient direction; ns = not significant. *** $p$<0.01, ** $p$<0.05 (robust standard errors). Salary excluded from all models (see Section~\ref{sec:notes}).
\end{tablenotes}
\end{threeparttable}
\end{table}


%─────────────────────────────────────────────────────────────────────────────
\section{Next Steps}
%─────────────────────────────────────────────────────────────────────────────

The associations documented here, particularly the divergent roles of social factors and
professional development across the wage floor and ceiling, provide a useful
starting point for designing a pilot RCT that could isolate the mechanisms at play.
Expanding the sample and cross-tabulating wage WTA by field of study and graduation
year are natural immediate next steps before pilot design.

\end{document}
